\section{Анализ предметной области}

\subsection{Интерактивные приложения как класс программных систем}

Интерактивным приложением называют программную систему, результат
работы которой зависит не только от заранее заданных данных, но и от
постоянного взаимодействия с пользователем. Такое приложение получает
события от клавиатуры, мыши или другого устройства ввода, изменяет
внутреннее состояние и выводит обновленное изображение с достаточно малой
задержкой, чтобы пользователь воспринимал происходящее как непрерывный
процесс.

К интерактивным приложениям относятся компьютерные игры, учебные
симуляторы, демонстрационные 3D-сцены, инженерные визуализаторы,
виртуальные лаборатории, средства прототипирования интерфейсов и
интерактивные художественные инсталляции. Эти программы различаются по
назначению, но имеют общую техническую основу. Для них характерны
периодическое обновление состояния, отрисовка кадра, обработка ввода и
управление набором объектов, присутствующих в сцене.

Обычное прикладное приложение часто строится вокруг запроса
пользователя: пользователь нажимает кнопку, программа выполняет действие
и возвращается в состояние ожидания. Интерактивное приложение работает
иначе. В нем присутствует непрерывный цикл: программа опрашивает
устройства ввода, рассчитывает изменения за прошедший промежуток времени,
обновляет положение и состояние объектов, после чего формирует новый кадр.
Такой цикл принято называть главным циклом приложения или игровым
циклом, даже если разрабатываемая программа не является игрой.

В простейшем виде работа интерактивного приложения может быть
представлена схемой, приведенной на рисунке~\ref{fig:interactive-loop}.

\begin{figure}[H]
\centering
\begin{tabular}{c}
\fbox{\parbox{0.78\textwidth}{\centering получение событий ввода}}\\[3mm]
$\downarrow$\\[3mm]
\fbox{\parbox{0.78\textwidth}{\centering обновление состояния объектов сцены}}\\[3mm]
$\downarrow$\\[3mm]
\fbox{\parbox{0.78\textwidth}{\centering обработка поведения, столкновений и зависимостей}}\\[3mm]
$\downarrow$\\[3mm]
\fbox{\parbox{0.78\textwidth}{\centering формирование и вывод кадра}}\\[3mm]
$\downarrow$\\[3mm]
\fbox{\parbox{0.78\textwidth}{\centering переход к следующей итерации цикла}}
\end{tabular}
\caption{Обобщенная схема работы интерактивного приложения}
\label{fig:interactive-loop}
\end{figure}

На практике каждый из показанных этапов распадается на несколько
подзадач. Обработка ввода включает регистрацию нажатий клавиш,
перемещения мыши, изменения состояния кнопок и колесика прокрутки.
Обновление состояния объектов связано с вычислением трансформаций,
анимаций, перемещений, пользовательских сценариев поведения. Обработка
столкновений требует определения пересечений между объектами, реакции на
контакт и, при необходимости, изменения скорости или положения объектов.
Отрисовка кадра предполагает подготовку данных для графического
процессора, выбор камеры, материалов, источников света, шейдеров и
очередности вывода объектов.

Особенность таких приложений состоит в том, что их внутренняя
структура быстро усложняется с ростом числа объектов. Один объект может
иметь положение в пространстве, графическое представление, коллайдер,
физическое тело, пользовательский сценарий поведения, источник света или
камеру. При этом разные объекты используют разные сочетания этих
возможностей. Например, камера не должна иметь графическую модель, а
декоративный объект может не участвовать в столкновениях. Если каждую
разновидность объекта описывать отдельным классом, то количество классов
и связей между ними начинает расти быстрее, чем сама предметная область.
Именно поэтому при разработке интерактивных приложений особенно важен
выбор архитектурного подхода.

\subsection{Программные средства построения интерактивных приложений}

Разработка интерактивного приложения с нуля требует решения большого
числа технических задач, не связанных напрямую с логикой конкретного
проекта. Разработчику нужно создать окно, получить графический контекст,
организовать цикл обновления, загрузить ресурсы, отрисовать геометрию,
обработать ввод, хранить состояние сцены, описать поведение объектов и
обеспечить освобождение ресурсов. Поэтому на практике редко используют
только стандартную библиотеку языка программирования. Обычно выбирают
один из классов программных средств: низкоуровневую графическую
библиотеку, фреймворк, игровой движок или специализированную среду
прототипирования.

Низкоуровневая графическая библиотека предоставляет доступ к
аппаратно-ускоренной графике, но почти не навязывает архитектуру
приложения. Примером является OpenGL. Khronos Group описывает OpenGL как
кроссплатформенный API для отрисовки 2D- и 3D-графики
\cite{khronos-opengl}. Такой API решает задачу взаимодействия с графическим
процессором, но не предоставляет готовой модели сцены, системы сущностей,
иерархии объектов или компонентной архитектуры. Программист получает
контроль над рендерингом, но вынужден самостоятельно проектировать
остальную инфраструктуру приложения.

Для языка Java важную роль играет LWJGL. Официальный сайт
характеризует LWJGL как библиотеку, предоставляющую Java-приложениям
кроссплатформенный доступ к нативным API, включая OpenGL, OpenAL и OpenCL
\cite{lwjgl}. Фактически LWJGL является мостом между Java и низкоуровневыми
средствами работы с графикой, звуком и окном. Это делает библиотеку
удобной основой для собственных движков и графических инструментов, но
она не является готовой системой построения сцен. Разработчику необходимо
самостоятельно реализовывать структуру объектов, управление ресурсами,
обновление состояния и пользовательские абстракции.

Фреймворк находится на более высоком уровне. Он не только
предоставляет доступ к графике и вводу, но и задает часть структуры
приложения. Одним из наиболее известных Java-фреймворков является
libGDX. В официальном описании libGDX определяется как кроссплатформенный
Java-фреймворк для разработки игр на основе OpenGL ES, поддерживающий
Windows, Linux, macOS, Android, браузер и iOS \cite{libgdx}. Его сильная
сторона заключается в переносимости и широком наборе готовых средств.
Однако libGDX оставляет разработчику свободу в выборе архитектуры
предметных объектов. Для небольших проектов это удобно, но при построении
системы из множества объектов разработчик все равно должен заранее
решить, как именно будут описываться сущности, компоненты и связи между
ними.

Игровой движок поднимает уровень абстракции еще выше. Он обычно
содержит модель сцены, подсистемы рендеринга, ввода, ресурсов, анимации,
физики и инструменты для разработки. Для Java таким решением является
jMonkeyEngine. На официальном сайте он описан как современный игровой
движок, написанный преимущественно на Java и ориентированный на
разработчиков, которым нужен контроль над кодом и возможность расширять
движок под собственный процесс разработки \cite{jmonkeyengine}. Это
полноценный инструмент, но вместе с полнотой возможностей появляется и
более высокий порог входа. Пользователь должен принять архитектурные
решения движка, изучить его жизненный цикл, систему ресурсов,
пространственную модель и набор классов.

Отдельную группу составляют среды прототипирования и визуального
программирования. Processing, согласно официальному описанию, является
гибким программным «скетчбуком» и языком для обучения программированию,
который с 2001 года используется в визуальных искусствах, образовании и
быстром прототипировании \cite{processing}. Такие инструменты снижают
барьер входа, но их модель обычно ориентирована на небольшие эскизы,
визуальные эксперименты и учебные проекты. Для разработки расширяемой
компонентной системы они подходят хуже, чем библиотека или фреймворк,
поскольку скрывают часть инфраструктуры и задают собственный стиль
организации программы.

Таким образом, между низкоуровневой библиотекой и полноценным игровым
движком существует промежуточная ниша. Java-разработчику может быть
нужна не готовая тяжелая среда с редактором и большим набором подсистем, а
легковесная программная основа, позволяющая строить интерактивные сцены
из сущностей, компонентов и систем. Такая основа должна закрывать типовые
задачи цикла приложения, сцены, рендеринга, ввода и базовой обработки
столкновений, но при этом оставаться понятной и расширяемой на уровне
исходного кода.

\subsection{Компонентно-ориентированный подход к построению интерактивных приложений}

Одним из главных вопросов при проектировании интерактивного
приложения является способ описания объектов сцены. В традиционном
объектно-ориентированном подходе разработчик часто начинает с иерархии
классов: базовый класс объекта, затем классы персонажей, снарядов,
декораций, источников света, камер, физических объектов. На раннем этапе
такая схема кажется удобной. Например, можно создать базовый класс
\texttt{GameObject}, от него унаследовать \texttt{RenderableObject}, затем
\texttt{PhysicalRenderableObject} и так далее.

Проблема появляется, когда свойства начинают комбинироваться. Объект
может быть видимым, но не иметь физического тела. Другой объект может
участвовать в столкновениях, но не отображаться на экране. Третий объект
может иметь скрипт поведения и источник света. При наследовании такие
сочетания приводят либо к очень глубокой иерархии, либо к дублированию
кода в разных ветвях. Изменение одного свойства становится изменением
класса, а не добавлением новой возможности к объекту.

Компонентный подход решает эту проблему иначе. Объект сцены
рассматривается не как экземпляр специализированного класса, а как
контейнер для набора независимых компонентов. Компонент описывает одну
характеристику или одну сторону поведения объекта: положение в
пространстве, камеру, отрисовку сетки, источник света, физическое тело,
коллайдер, пользовательский скрипт. Чтобы изменить объект, разработчик
добавляет или удаляет компоненты, не создавая новый класс для каждого
сочетания возможностей.

Преимущество компонентной модели хорошо видно на простом примере.
Объект, имеющий компоненты \texttt{Transform} и \texttt{MeshRenderer},
может быть отрисован как обычная модель. Если к нему добавить компонент
\texttt{Rigidbody}, он начинает участвовать в физической обработке. Если
добавить компонент \texttt{Script}, объект получает пользовательское
поведение. При этом базовая сущность остается той же самой, а набор
возможностей определяется составом компонентов.

В популярных инструментах разработки интерактивных приложений
компонентная идея встречается достаточно часто. Например, в Unity
функциональность объекта сцены определяется компонентами, прикрепленными
к GameObject \cite{unity-components}. Unity также развивает отдельное
направление ECS, называя его дата-ориентированным фреймворком для
высокопроизводительной обработки сущностей \cite{unity-ecs}. Однако Unity
не является Java-решением и представляет собой крупную экосистему с
собственным редактором, форматом проектов и моделью разработки. Для
Java-разработчика, желающего изучить или использовать компонентную модель
в небольшом приложении, такой инструмент может оказаться избыточным.

Компонентно-ориентированную архитектуру необходимо отличать от
канонического Entity-Component-System. В строгом варианте ECS сущность
обычно является только идентификатором, компонент хранит данные, а вся
логика выполняется системами, которые выбирают сущности по набору
компонентов. Такой подход хорошо подходит для высокопроизводительных
сцен, поскольку данные можно располагать в памяти компактно и обрабатывать
пакетами. Но он требует более строгой организации данных и усложняет
разработку для небольших проектов.

Разрабатываемая программная система ближе к компонентно-ориентированной
модели с элементами ECS, чем к полностью каноническому
дата-ориентированному ECS. В ней есть сущности, компоненты, системы и
запросы к сущностям по составу компонентов. При этом компоненты являются
объектами Java и могут иметь собственный жизненный цикл. Такой вариант
сохраняет основную идею композиции объектов из компонентов, но остается
более привычным для Java-разработчика, работающего с объектами, классами
и методами.

Для ВКР выбран именно этот подход, так как он позволяет совместить
простоту объектной модели Java с гибкостью компонентной структуры сцены.
Разработчик получает возможность описывать интерактивные объекты через
состав компонентов, не создавая сложную иерархию наследования и не
переходя к полностью дата-ориентированной архитектуре.

\subsection{Ключевые сущности предметной области}

Несмотря на различие конкретных библиотек и движков, большинство
средств построения интерактивных приложений оперирует схожими
понятиями. Ключевой сущностью является сцена. Она задает контекст
выполнения: содержит набор объектов, управляет их жизненным циклом,
обновляет системы и определяет, какие объекты участвуют в текущем кадре.
В играх сцена может соответствовать уровню, меню или отдельному режиму.
В учебном симуляторе сцена может представлять лабораторную установку,
эксперимент или демонстрационный пример.

Сущность представляет отдельный объект сцены. Сама по себе сущность
может не содержать логики и данных предметной области. Ее назначение
состоит в том, чтобы объединять компоненты и выступать единицей
идентификации. Такой подход позволяет рассматривать любой объект сцены
одинаково: камеру, куб, источник света, физическое тело или служебный
объект управления.

Компонент описывает отдельную характеристику сущности. Наиболее
универсальным компонентом является трансформация, задающая положение,
поворот и масштаб объекта. Без нее невозможно определить, где находится
модель, камера или коллайдер. Другие компоненты добавляют графическое
представление, камеру, освещение, физическое тело, обработчик столкновений
или пользовательское поведение.

Система выполняет обработку сущностей, имеющих определенный набор
компонентов. Например, система рендеринга выбирает сущности с
трансформацией и графическим компонентом, система скриптов вызывает
пользовательское поведение, а физическая система работает с физическими
телами и коллайдерами. Благодаря этому логика обработки не размазывается
по всем объектам, а сосредотачивается в подсистемах, отвечающих за
конкретную область.

Ресурсы приложения включают шейдеры, текстуры, материалы, меши и другие
данные, используемые несколькими объектами. В интерактивных приложениях
важно не только загрузить ресурс, но и контролировать срок его жизни.
Графические ресурсы занимают память не только в оперативной памяти
приложения, но и на стороне графического процессора, поэтому при смене
сцены или завершении приложения их необходимо корректно освобождать.

Камера определяет способ наблюдения сцены. В 3D-приложении она хранит
параметры проекции, область видимости, соотношение сторон и положение в
пространстве. Источники света влияют на расчет освещенности объектов.
Материал и шейдер задают способ визуализации поверхности. Ввод
пользователя связывает внешние действия с изменениями состояния сцены.

Связь этих сущностей можно описать так: приложение содержит активную
сцену, сцена содержит сущности и системы, сущности состоят из компонентов,
а системы обрабатывают подходящие сущности и используют ресурсы для
формирования результата. Такая модель достаточно проста для понимания, но
при этом позволяет строить разнообразные интерактивные приложения без
создания отдельного класса для каждого вида объекта.

\subsection{Особенности разработки интерактивных приложений на Java}

Java традиционно используется для разработки серверных приложений,
корпоративных систем, инструментов автоматизации, учебных программ и
кроссплатформенных настольных приложений. Для интерактивной графики Java
применяется реже, чем C++ или C\#, однако у нее есть ряд свойств, полезных
для рассматриваемой задачи. Язык имеет строгую типизацию, развитую
стандартную библиотеку, автоматическое управление памятью и широкий набор
средств разработки. Официальная документация Java SE 25 включает
спецификации платформы, API и инструменты JDK, что делает платформу
достаточно полно документированной для учебного и прикладного
использования \cite{oracle-java25}.

Главное преимущество Java в контексте данной работы заключается в
понятной объектной модели. Сущности, компоненты, системы, материалы,
меши, камеры и ресурсы естественно выражаются через классы и интерфейсы.
Это упрощает проектирование программной системы, ориентированной на
разработчика, который уже знаком с Java и хочет создавать интерактивные
сцены без перехода на другой язык и без изучения крупной игровой среды.

При этом Java не предоставляет встроенного доступа к OpenGL на уровне
стандартной библиотеки. Для работы с аппаратно-ускоренной графикой нужны
дополнительные библиотеки, такие как LWJGL. Использование LWJGL позволяет
сохранить разработку на Java, но работать с графическим API, близким к
нативному уровню. Такой подход требует большей аккуратности: разработчик
должен учитывать создание окна, контекста OpenGL, буферов, шейдеров,
загрузку данных на графический процессор и освобождение ресурсов.

Еще одна особенность связана с производительностью. Интерактивное
приложение должно обновлять состояние и отрисовывать кадры много раз в
секунду. Поэтому неудачные архитектурные решения быстро становятся
заметными: лишние выделения памяти, частый поиск объектов, неоптимальная
сортировка команд рендеринга или избыточная работа с графическими
ресурсами приводят к задержкам. В Java часть низкоуровневых деталей
скрыта виртуальной машиной, но это не отменяет необходимости аккуратно
строить цикл обновления и ограничивать лишние операции внутри него.

Для рассматриваемой программной системы важно, что математика
трехмерной сцены реализуется в самом проекте. Векторные, матричные и
кватернионные операции используются для трансформаций, камеры,
ориентации объектов и расчета положения в пространстве. Такой подход
увеличивает объем реализации, но делает систему более самостоятельной и
позволяет подробно раскрыть проектные решения в техническом проекте и
рабочем проекте.

Отсутствие готового сборщика на раннем этапе разработки не влияет на
анализ предметной области, но отражает текущий уровень зрелости
программной системы. Для дальнейшего развития проекта целесообразно
добавить Maven или Gradle, чтобы упростить подключение зависимостей,
запуск тестов и воспроизводимую сборку. В рамках первой главы этот вопрос
важен как подтверждение общей проблемы: Java-разработчику нужны не
только низкоуровневые библиотеки, но и удобная, воспроизводимая
инфраструктура создания интерактивных приложений.

\subsection{Анализ существующих программных решений}

Для сравнения выбраны решения, близкие к теме работы по разным
признакам. LWJGL представляет низкоуровневый доступ к графическим и
системным API из Java. libGDX является Java-фреймворком для
кроссплатформенной разработки. jMonkeyEngine представляет полноценный
Java-движок. Processing отражает подход быстрого визуального
прототипирования на Java-подобной основе. Unity включен в анализ не как
Java-решение, а как распространенный пример компонентной модели и
практики построения интерактивных сцен.

Краткое сравнение программных средств построения интерактивных
приложений приведено в таблице~\ref{tab:tools-comparison}.

\begingroup
\small
\setlength{\tabcolsep}{4pt}
\renewcommand{\arraystretch}{1.2}

\begin{longtable}{|
>{\raggedright\arraybackslash}p{0.14\textwidth}|
>{\raggedright\arraybackslash}p{0.23\textwidth}|
>{\raggedright\arraybackslash}p{0.25\textwidth}|
>{\raggedright\arraybackslash}p{0.28\textwidth}|}
\caption{Сравнение программных средств построения интерактивных приложений}
\label{tab:tools-comparison}\\
\hline
\centrow Средство &
\centrow Тип и платформа &
\centrow Возможности &
\centrow Особенности для Java-разработчика \\
\hline
\endfirsthead

\multicolumn{4}{l}{Продолжение таблицы~\ref{tab:tools-comparison}}\\
\hline
\centrow Средство &
\centrow Тип и платформа &
\centrow Возможности &
\centrow Особенности для Java-разработчика \\
\hline
\endhead

\hline
\endlastfoot

LWJGL &
Низкоуровневая Java-библиотека, использующая нативные привязки &
Доступ к OpenGL, OpenAL и другим системным API; построение 2D- и 3D-графики возможно через OpenGL; собственного редактора нет &
Дает высокий уровень контроля над графикой и окружением выполнения, но не предоставляет готовую сцену, компонентную модель, управление ресурсами и жизненный цикл объектов. \\ \hline

libGDX &
Java-фреймворк для разработки игр и графических приложений &
Поддерживает 2D- и 3D-графику, работу с ресурсами, вводом и несколькими целевыми платформами; обязательного редактора нет &
Удобен для кроссплатформенной разработки, однако архитектура игровых объектов и способ организации логики во многом остаются ответственностью разработчика. \\ \hline

jMonkeyEngine &
Java-движок для создания трехмерных приложений &
Ориентирован преимущественно на 3D-графику, содержит готовые подсистемы сцены, материалов, освещения, физики и ресурсов; имеются дополнительные инструменты разработки &
Близок к полноценному решению на Java, но требует изучения крупного фреймворка и принятия его проектной модели, что может быть избыточно для небольших учебных и исследовательских проектов. \\ \hline

Processing &
Среда и библиотека для визуального программирования на Java-подобной основе &
Позволяет быстро создавать 2D- и 3D-скетчи, визуализации и учебные интерактивные примеры; имеет собственную среду разработки &
Имеет низкий порог входа, но больше ориентирован на визуальные эксперименты и прототипирование, чем на построение расширяемой компонентной архитектуры. \\ \hline

Unity &
Крупный игровой движок и редактор, использующий язык C\# &
Поддерживает 2D- и 3D-графику, физику, визуальный редактор сцен, компонентную модель объектов и развитую экосистему инструментов &
Показывает преимущества компонентного подхода, но не относится к Java-экосистеме и может быть слишком тяжелым решением для небольших проектов, где требуется простая программная основа. \\ \hline

\end{longtable}

\endgroup

Из сравнения видно, что существующие решения закрывают разные части
задачи. LWJGL удобен как низкоуровневая основа, но не решает архитектурные
вопросы. libGDX снижает объем технической работы и дает готовый фреймворк,
но не ставит компонентную модель в центр архитектуры. jMonkeyEngine
наиболее близок к полноценному движку на Java, однако его использование
предполагает освоение более крупной системы. Processing удобен для быстрых
визуальных экспериментов, но не ориентирован на построение расширяемой
сценной архитектуры. Unity показывает преимущества компонентного подхода,
но относится к другой технологической экосистеме.

Для небольшого Java-проекта возникает разрыв между двумя крайностями.
С одной стороны, можно взять низкоуровневую библиотеку и написать всю
архитектуру самостоятельно. С другой стороны, можно использовать
полноценный движок, приняв его сложность и набор инструментов. В рамках
ВКР интерес представляет промежуточное решение: программная система,
которая предоставляет компонентную организацию сцены, базовые системы
рендеринга, ввода, скриптов и обработки столкновений, но остается
достаточно легкой для изучения и модификации.

\subsection{Выводы по анализу предметной области}

Интерактивные приложения отличаются от обычных прикладных программ
наличием постоянного цикла обработки ввода, обновления состояния и вывода
кадра. Для их разработки требуется не только реализовать предметную логику,
но и создать инфраструктуру сцены, ресурсов, объектов, компонентов,
рендеринга и пользовательского ввода.

Анализ существующих решений показал, что Java-разработчик может
использовать низкоуровневые библиотеки, фреймворки или полноценные
движки. Однако низкоуровневые библиотеки требуют самостоятельной
разработки архитектуры, а крупные движки часто оказываются избыточными
для небольших учебных, демонстрационных и исследовательских проектов.
Это создает потребность в промежуточной программной системе, которая
закрывает основные инфраструктурные задачи, но остается простой для
изучения и расширения.

Компонентно-ориентированный подход позволяет отказаться от сложных
иерархий наследования и описывать объекты сцены как набор независимых
компонентов. Для выбранной темы это особенно важно, поскольку объекты
интерактивной сцены отличаются не типом класса, а набором возможностей:
одни отображаются, другие участвуют в столкновениях, третьи выполняют
пользовательскую логику или задают камеру.

Разрабатываемая программная система должна быть ориентирована на
Java-разработчика и предоставлять базовый набор средств для построения
интерактивных сцен: управление сценой, сущностями, компонентами и
системами, рендеринг, обработку ввода, скрипты поведения и базовую
обработку столкновений. Такой результат соответствует утвержденной теме
ВКР и создает основу для дальнейшего развития системы до более
полноценного игрового или интерактивного движка.
